\setcounter{chapter}{10}
\chapter{Les lois de Képler}

\section{Un satellite de la Terre}
\subsection{Situation}

On étudie le système Satellite (noté S) de masse $m$ situé à une distance $r$ du centre de la Terre (noté T) et de masse $M_T$. Le référentiel choisit est le référentiel: \trou{géocentrique} considéré galiléen.


\blocrouge{Définition}{le référentiel géocentrique est situé au centre de la Terre et possède trois axes fixes dans l'espace. Il est galiléen.
}

\paragraph*{Bilan des forces} Le système est soumis à la force de gravitation exercée par la Terre sur le satellite. 

\begin{figure}[!htp]
\begin{center}
\includegraphics[height=5cm]{gfx/gravitation}
\end{center}
\caption{Force d'attraction gravitationnelle}
\end{figure}
\subsection{Expression vectorielle de la force gravitationnelle}

On sait que la valeur de la force gravitationnelle subit par le Satellite vaut 
\begin{equation*}
	\boxed{F_{Terre/Satellite}=\trou{\mathcal{G}.\frac{m\times M_T}{r^2}} } 
\end{equation*}

\begin{itemize}
	\item m, $M_T$ les masses des objets en kg
	\item $r$: la distance en m entre les centres des objets
	\item $\mathcal{G}=\SI{6,67e-11}{S.I}$: la constante universelle de gravitation à ne pas confondre avec $g$ l'intensité du champ de pesanteur.
 
\end{itemize}

On doit être capable d'écrire mathématiquement $\overrightarrow{F_{Terre/Satellite}}$. Pour cela on utilise un vecteur unitaire (de norme \trou{égale à 1})que l'on indique sur le schéma. Ici, c'est le vecteur $\overrightarrow{u_{ST}}$  orienté de \trou{ S vers T}. On obtient alors:

$$\overrightarrow{F_{Terre/Satellite}}=\trou{\mathcal{G}.\frac{mM_T}{r^2}\overrightarrow{u_{ST}}}$$
% \marginpar{Cette expression est très souvent demandée en exercice}
\subsection{Expression vectorielle de l'accélération}
Appliquons la seconde de Newton au système satellite:
\begin{eqnarray}
m \overrightarrow{a} && = \overrightarrow{F_{Terre/Satellite}} \nonumber \\
&& = \mathcal{G}.\frac{m.M_T}{r^2}. \overrightarrow{u_{ST}} \nonumber
\end{eqnarray}
En simplifiant par $m$ on trouve l'expression vectorielle du vecteur accélération:


\begin{equation} \overrightarrow{a} =\trou{\mathcal{G}.\frac{M_T}{r^2}.\overrightarrow{u_{ST}}} \label{eq:acc}
\end{equation}

\subsection{Calcul de la vitesse}
\paragraph*{Rappel de cinématique}
Pour un mouvement circulaire quelconque le vecteur accélération $ \overrightarrow{a}$ se décompose en une composante tangentielle $a_T$ et une composante normale $ a_N$ telles que: 
\begin{figure}[!h]
\begin{center}
\includegraphics[scale=0.7]{gfx/mvtcirculaire}
\end{center}
\caption{Composantes du vecteur accélération lors d'un mouvement circulaire}
\end{figure}

\begin{equation}
\overrightarrow{a} =\trou{ \frac{dv}{dt}}. \overrightarrow{T} + \trou{\frac{v^2}{R}} . \overrightarrow{N}
\label{eq:frenet}
\end{equation}

$ \overrightarrow{T}$ et $ \overrightarrow{N}$ sont des vecteurs unitaires (de norme égale à 1) ils sont appelés vecteur Tangentiel et vecteur Normal.

 En comparant les équations \ref{eq:acc} et \ref{eq:frenet} et en constatant que $ \overrightarrow{N}= \overrightarrow{u_{ST}}$, on trouve par identification:
 \begin{eqnarray}
 a_T && = \frac{dv}{dt}= \trou{0} \label{eq:aT}\\
 a_N && = \frac{v^2}{r}=\trou{ \mathcal{G}\frac{M_T}{r^2}} \label{eq:vitesse}
 \end{eqnarray}
De l'équation \ref{eq:aT}, on en déduit que \trou{ $v=cste$} c'est à dire que le mouvement est \trou{uniforme}.
Nous venons de prouver que si le mouvement est \trou{circulaire} alors il est \trou{uniforme}.
%\marginpar {Cette preuve est à savoir faire pour le bac}


A l'aide de l'équation \ref{eq:vitesse}, on isole $v$ et on en déduit la vitesse d'un satellite en mouvement circulaire uniforme dont le rayon de l'orbite est $r$:
\begin{equation}
\boxed{v=\sqrt{\mathcal{G}\frac{M_T}{r}}}
\end{equation}
\etiquette{Conséquence} Plus le satellite est éloigné plus sa vitesse est \trou{petite} (r au dénominateur). Plus l'astre autour duquel il tourne est massif et plus sa vitesse est \trou{grande} ($M_T$ au numérateur)
\subsection{Calcul de la période}
Le mouvement étant circulaire uniforme, S effectue un mouvement périodique autour de la Terre. 

\blocrouge{Définition}{  La durée notée $T$ mise par S pour faire une fois le tour de la Terre, s'appelle la \trou{\textbf{période de révolution}}.}

 
 Puisque la distance parcourue autour de la Terre vaut: $2\pi r$ et que cette distance a été parcourue à une vitesse $v$ pendant une durée $T$, on a:
 
\begin{eqnarray}
T  && =\trou{ \frac{2 \pi r}{v} }\label{eq:periode}\\
T  && = 2\pi \sqrt{\frac{r^3}{\mathcal{G}.M_T}} \label{eq:periodeG}
\end{eqnarray}

\newpage
\blocgris{Application}{On considère l'orbite de la Lune comme étant circulaire de rayon $r=\SI{3.84e5}{\kilo\meter}$. La vitesse de la Lune autour de la Terre est : $v_L=\SI{1.02e3}{\meter\per \second}$

\begin{enumerate}
\item Exprimer, puis calculer la période de révolution $T_L$ de la Lune autour de la Terre. Convertir cette durée en jours.
\item La période de rotation de la Lune sur elle-même est égale à 27,3 jours. Comparer cette période à $T_L$. Quelle est la conséquence de ce résultat ?
\end{enumerate}
}
\vspace{6cm}
\section{Lois de Képler}
\subsection{Référentiel héliocentrique}
Pour décrire le mouvement des planètes autour du Soleil, il est pratique d'affirmer que le Soleil est immobile est que les planètes sont en mouvement autour de lui. Cette description revient à choisir comme référentiel le Soleil, on parle de référentiel \trou{héliocentrique}. Les axes de ce référentiel pointent vers 3 étoiles éloignées (et donc considérées comme fixes). On admettra que ce référentiel est galiléen.

Simulateurs des lois de Képler: 

\url{https://phet.colorado.edu/sims/html/keplers-laws/latest/keplers-laws_all.html}
\\ \url{https://physique.ostralo.net/kepler/}
\subsection{Première loi (1609)}

\blocrouge{Première Loi de Képler}{
Dans le référentiel héliocentrique, la trajectoire du centre de gravité d'une planète est \trou{une ellipse} dont le centre du Soleil est un \trou{Foyer}}
%\marginpar{Les lois de Képler sont à savoir}
\newpage
\subsection{Seconde loi}

\begin{figure}[!htp]
\begin{center}
\includegraphics[height=4cm]{gfx/Loidesaires}

\end{center}
\caption{2nde loi de Képler}
\label{fig:loidesaires}
\end{figure}

\blocrouge{Seconde Loi de Képler}{
Le segment de droite reliant les centres de gravité du Soleil et de la planète balaie \trou{des aires égales pendant des durées égales}.} 


\etiquette{Explications}

 Soit $\Delta t$ une durée (par exemple 1 mois), on a représenté sur la figure \ref{fig:loidesaires} l'aire balayée par le segment reliant le centre de la Terre et le centre du Soleil. Ce segment a décrit une surface\trou{ $A_1$}. 

D'après la 2nde loi de Képler l'aire décrite par ce segment à partir d'un autre instant $t_2$ et pour une même durée $\Delta t$ vaut \trou{$A_2$ } et on a \trou{$A_1=A_2$}.

Or, la distance entre B et C étant plus grande qu'entre D et E cela signifie que la vitesse de la Terre est plus \trou{grande} sur la portion BC que sur la portion DE. Dans ce cas, on prouve que lorsque le mouvement est elliptique, il n'est pas \trou{uniforme}.


\etiquetterouge{Conclusion} La deuxième loi de Képler permet de justifier que la Terre n'a pas toujours la même vitesse au cours de son mouvement autour du Soleil. Sa vitesse est plus \trou{grande} lorsqu'elle est \trou{proche} du Soleil.

\subsection{Troisième loi de Képler(1618)}

\blocrouge{Enoncé }{
Pour toutes les planètes gravitant autour du Soleil à la période T et dont le demi grand axe de la trajectoire vaut a, on a: $$\trou{\frac{T^2}{a^3}=constante}$$
}
\etiquette{Cas particulier important: MCU}

Dans le cas du MCU, le rayon du cercle joue le rôle du demi-grand axe $a$.
On peut alors démontrer la troisième loi de Képler dans le cas particulier de la trajectoire circulaire. Il suffit de mettre au carré l'équation  \ref{eq:periodeG}. 

\begin{eqnarray}
T^2 && =4\pi^2\frac{r^3}{\mathcal{G}.M_T}  \\
\frac{T^2}{r^3} && =\frac{4 \pi^2}{\mathcal{G}.M_T}=cste 
\end{eqnarray}

\etiquetterouge{Les points importants}

Les choses à connaître par coeur:
\begin{itemize}
	\item Les définitions
	\item L'expression vectorielle de la force de gravitation
	\item Les coordonnées du vecteur accélération pour un mouvement circulaire uniforme dans le repère de Frénet
	\item Les 3 lois de Képler + schémas
\end{itemize}

\etiquetterouge{A savoir faire}

\begin{itemize}
	\item Ecrire l'expression vectorielle de la force de gravitation en prenant en compte le sens du vecteur unitaire choisi ou imposé par le texte.
	\item Démontrer que si le mouvement est circulaire alors la vitesse de l'astre (ou satellite) est constante.
	\item Démontrer la troisième loi de Képler dans le cas d'une trajectoire circulaire
\end{itemize}

\blocgris{Sujets de bac pour s'entraîner}{

\begin{itemize}
\item  Polynésie Jour 2 2024 "La masse de la Terre" 
	\item Bac 2024 Centres étrangers 1 Jour 2 EXERCICE 2 : à la découverte des lunes glacées de Jupiter (6 points)
\item Asie 2014
Exercice 1 Le monde selon Hubble (7 points)
\end{itemize}
}